Consideriamo un mezzo poroso saturo, deformabile e a piccole deformazioni,
nel quale il fluido interstiziale trasporta un soluto che non retroagisce
sulla risposta idromeccanica. La formulazione combina la consolidazione
quasi-statica di Biot~\citep{biot1941general,bociu2023mathematical}, una
legge costitutiva di Kelvin--Voigt con memoria di Caputo per lo scheletro
solido~\citep{caponi2021sapio,failla2020advanced} e un'equazione di
trasporto con accumulo temporale frazionario. La scelta modellistica
consiste nel collocare le due memorie in \emph{sottosistemi distinti},
lasciando locale nel tempo la legge di Darcy. In particolare, la
formulazione non introduce una memoria ulteriore nella permeabilit\`a o
nell'equazione di accumulo del fluido.

\subsection{Ipotesi, derivate frazionarie e legge costitutiva}
\label{sec:forward-constitutive}

Siano $\Omega\subset\R^m$, $m\in\{2,3\}$, un dominio lipschitziano limitato
e $T>0$. Per $(\mathbf x,t)\in\Omega\times(0,T]$, le incognite sono lo
spostamento $\mathbf d$, la pressione interstiziale $p$ e la concentrazione
nel fluido $c$; il flusso di Darcy $\mathbf q$ \`e ricostruito dalla pressione.
La cinematica \`e linearizzata e il tensore delle piccole deformazioni \`e
\begin{equation}
 \eps(\mathbf d)=\tfrac12(\nabla\mathbf d+\nabla\mathbf d^{\mathsf T}).
 \label{eq:strain}
\end{equation}
La geometria di riferimento resta fissa: porosit\`a e coefficienti
materiali sono campi prescritti, indipendenti dal tempo, e non vengono
aggiornati mediante una legge di evoluzione della porosit\`a. Questa
approssimazione non rappresenta una teoria a deformazioni finite della
miscela solido--fluido.

Per una funzione assolutamente continua $u:[0,T]\to X$, a valori in uno
spazio di Hilbert $X$, poniamo, per $0<\gamma<1$,
\begin{equation}
 \CD{\gamma}u(t)=\frac{1}{\Gamma(1-\gamma)}
 \int_0^t(t-s)^{-\gamma}\,\partial_su(s)\,\dd s
 = (k_\gamma*\partial_tu)(t),
 \label{eq:caputo}
\end{equation}
con $k_\gamma(t)=t^{-\gamma}/\Gamma(1-\gamma)$ per $t>0$.
Le derivate agiscono componente per componente; per $\gamma=1$ si
pone $\CD1 u=\partial_tu$. Ne segue che $\CD\gamma u=0$ per una storia
costante e che la memoria dipende dall'intera storia successiva a $t=0$.
Sotto la regolarit\`a assunta, derivata temporale, operatore di
deformazione e moltiplicazione per coefficienti indipendenti dal tempo
commutano.

La tensione totale \`e scomposta in una parte efficace e nella pressione
di Biot secondo
\begin{equation}
 \sig=\sig_{\rm eff}-bp\I,\qquad
 \sig_{\rm eff}
 =E(\mathbf x)\mathbb C_0:\eps(\mathbf d)
 +A_v\mathbb C_0:\CD\beta\eps(\mathbf d).
 \label{eq:dim-visco}
\end{equation}
Qui $0<\beta\le1$ governa la memoria dello scheletro,
$b\in(0,1]$ \`e il coefficiente di Biot, e
$A_v=\eta_v\tau_v^\beta>0$ \`e l'ampiezza viscoelastica, con
$\eta_v,\tau_v>0$. Se $\beta=1$, il secondo addendo \`e il termine
viscoso di Kelvin--Voigt ordinario; per $0<\beta<1$ \`e una
convoluzione con un nucleo a legge di potenza. La formulazione non
comporta un modulo di rilassamento uniforme in $\beta$, poich\'e anche
$A_v$ dipende da tale ordine quando $\eta_v$ e $\tau_v$ sono fissati.

Il tensore isotropo adimensionale $\mathbb C_0$ soddisfa, su ogni tensore
simmetrico $e$,
\begin{equation}
 \mathbb C_0:e=\frac{e}{1+\nu}
 +\frac{\nu\,\mathrm{tr}(e)}{(1+\nu)(1-2\nu)}\I,
 \qquad -1<\nu<\tfrac12 .
 \label{eq:C0}
\end{equation}
La scelta $\nu\in(-1,1/2)$ assicura la coercivit\`a della legge elastica
isotropa tridimensionale; nel caso $m=2$ si impone \emph{deformazione piana},
ossia si usa la restrizione della legge 3D con deformazione fuori piano
nulla. L'ampiezza $A_v$ moltiplica lo stesso tensore $\mathbb C_0$,
ma, diversamente dal modulo elastico, \`e uniforme nello spazio.

I coefficienti nominali positivi $E,K_h,D,\phi_0$ e i moltiplicatori
regionali $r_E,r_h,r_D,r_\phi$ sono distinti:
\begin{equation}
 \begin{aligned}
 E(\mathbf x)&=E\,r_E(\mathbf x),&
 \mathbb K_h(\mathbf x)&=K_h r_h(\mathbf x)\I,\\
 D(\mathbf x)&=D\,r_D(\mathbf x),&
 \phi(\mathbf x)&=\phi_0 r_\phi(\mathbf x).
 \end{aligned}
 \label{eq:material-coeffs}
\end{equation}
Assumiamo $r_j\in L^\infty(\Omega)$, con limiti inferiori
strettamente positivi uniformi. I campi possono essere discontinui alle
interfacce regionali. Nel caso omogeneo $r_j\equiv1$. In particolare,
la variazione regionale di $r_E$ non viene applicata al termine
frazionario: \`e una precisa ipotesi costitutiva, coerente con gli
operatori assemblati nel solutore.

\subsection{Sistema Biot--Darcy e trasporto del tracciante}
\label{sec:forward-system}

In variabili adimensionali, la formulazione forte \`e
\begin{subequations}\label{eq:coupled-system}
\begin{empheq}[left=\empheqlbrace]{align}
 -\nabla\!\cdot\sig &= \mathbf f_s, \label{eq:mech}\\
 \sig &= E(\mathbf x)\mathbb C_0:\eps(\mathbf d)-bp\I
       \notag\\[-2pt]
 &\quad +A_v\mathbb C_0:\CD\beta\eps(\mathbf d),
       \label{eq:stress}\\
 \partial_t(S_p p+b\nabla\!\cdot\mathbf d)+\nabla\!\cdot\mathbf q
       &=s_f, \label{eq:fluid}\\
 \mathbf q&=-\mathbb K_h(\mathbf x)\nabla p,
       \label{eq:darcy}\\
 \phi(\mathbf x)\CD\alpha c+\nabla\!\cdot\mathbf F_c
       &=s_c, \label{eq:tracer}
\end{empheq}
\end{subequations}
con $S_p>0$ coefficiente di immagazzinamento, $0<\alpha\le1$ ordine
frazionario del trasporto, $\mathbf f_s$ forza di volume e $s_f,s_c$
sorgenti prescritte. Il flusso totale del tracciante \`e
\begin{equation}
 \mathbf F_c=c\mathbf q-\phi(\mathbf x)D(\mathbf x)\nabla c .
 \label{eq:tracer-flux}
\end{equation}
Si mantiene la divergenza \emph{conservativa} dell'avvezione:
$\nabla\!\cdot(c\mathbf q)=\mathbf q\cdot\nabla c+
 c\nabla\!\cdot\mathbf q$ nei punti di regolarit\`a. La seconda
componente pu\`o essere non nulla durante la consolidazione;
ometterla modificherebbe il modello. La forma in divergenza definisce
anche il flusso normale attraverso le interfacce dei materiali.

L'operatore $\CD\alpha$ compare nell'accumulo, mentre avvezione e
diffusione rimangono locali nel tempo. Per $\alpha<1$,
\eqref{eq:tracer} \`e una \emph{chiusura fenomenologica a memoria}, non
la legge di continuit\`a istantanea $\partial_t(\phi c)+
\nabla\!\cdot\mathbf F_c=s_c$: le due equazioni coincidono soltanto
per $\alpha=1$. Con $\phi$ prescritta nel tempo, essa conserva
globalmente $\int_\Omega\phi c$ in assenza di sorgenti e di flusso
normale, come dimostrato pi\`u avanti. Per \emph{tracciante passivo}
intendiamo inoltre che $c$ non modifica $\sig$, $p$ o $\mathbf q$;
non viene modellata n\'e una cinetica di legame n\'e l'attraversamento
di una barriera vascolare.

L'accoppiamento \`e pertanto unidirezionale verso il trasporto:
per valori assegnati di $(E,\beta,K_h,A_v)$ e di carichi, il
sottosistema~\eqref{eq:mech}--\eqref{eq:darcy} determina
$(\mathbf d,p,\mathbf q)$ senza dipendere da $\alpha$ e $D$.
La successiva soluzione di~\eqref{eq:tracer} dipende da entrambi i
sottosistemi tramite $\mathbf q$. Se il problema idromeccanico
ammette soluzione unica, cambiare $\alpha$ a $\beta$ fissato lascia
inalterati spostamento, pressione e flusso. Non vale l'implicazione
inversa: un cambiamento di $\beta$ pu\`o risultare non osservabile
nella sola concentrazione, soprattutto se il flusso non trasmette
informazione ai sensori, come analizzato nella Sezione~\ref{sec:inverse}.

Questa separazione motiva quattro configurazioni a parametri restanti
fissati: M1 $(\alpha,\beta)=(1,1)$, M2 $(0.70,1)$,
M3 $(1,0.45)$ e M4 $(0.70,0.45)$. Nei confronti
numerici si assume $\tau_v=1$, cos\`i che $A_v=\eta_v$ rimane fissato quando si varia $\beta$. M1 \`e il modello
poro-viscoelastico con Kelvin--Voigt e trasporto ordinari,
M2 introduce la sola memoria del trasporto, M3 la sola memoria
meccanica frazionaria e M4 entrambe. Il confronto fattoriale
isola gli effetti dei singoli cambiamenti d'ordine; l'eventuale
non additivit\`a nella concentrazione \`e misurata sullo stesso
protocollo osservativo nella Sezione~\ref{sec:numerics}.

\subsection{Dimensioni fisiche, dati e condizioni al contorno}
\label{sec:forward-data}

Siano $L_0,T_0,E_0,c_{\rm ref}>0$ le scale di lunghezza, tempo,
tensione e concentrazione. Poniamo
$\mathbf x=L_0\mathbf x^*$, $t=T_0t^*$,
$\mathbf d=L_0\mathbf d^*$, $p=E_0p^*$,
$\mathbf q=(L_0/T_0)\mathbf q^*$ e
$c=c_{\rm ref}c^*$. Le stelle sono soppresse nelle
\eqref{eq:coupled-system}. Nel sistema dimensionale, il contributo
viscoelastico contiene
$A_{v,\rm phys}\CD\beta\eps(\mathbf d)$ con
$A_{v,\rm phys}=\eta_{v,\rm phys}\tau_{v,\rm phys}^\beta$,
$[\eta_{v,\rm phys}]=\mathrm{Pa}$ e
$[A_{v,\rm phys}]=\mathrm{Pa}\,\mathrm{s}^\beta$.
Analogamente, nell'accumulo dimensionale del tracciante
si adotta
\begin{equation}
 \phi\,\tau_c^{\alpha-1}\CD\alpha c
 +\nabla\!\cdot\mathbf F_c=s_c,
 \qquad [\tau_c]=\mathrm{s} .
 \label{eq:dim-tracer}
\end{equation}
Il prefattore $\tau_c^{\alpha-1}$ riporta l'accumulo alle
unit\`a di concentrazione per tempo. Dividendo l'intera equazione per tale coefficiente, anche i
termini di flusso e sorgente vengono riscalati: il coefficiente
diffusivo cos\`i ottenuto ha unit\`a $L^2/T^\alpha$ e vale
$D_{\rm phys}\tau_c^{1-\alpha}$. Non si identifica in
questo modo una seconda diffusivit\`a indipendente. La scelta modellistica
$T_0=\tau_c$ conduce alle conversioni
\begin{equation}
\begin{aligned}
 E&=E_{\rm phys}/E_0,&
 S_p&=E_0 S_{p,\rm phys},\\
 D&=D_{\rm phys}T_0/L_0^2,&
 \mathbb K_h&=\frac{E_0T_0}{L_0^2}\mathbb K_{h,\rm phys},\\
 \eta_v&=\eta_{v,\rm phys}/E_0,&
 \tau_v&=\tau_{v,\rm phys}/T_0.
\end{aligned}
\label{eq:scales}
\end{equation}
Per il termine dimensionale di forza, per la sorgente di fluido,
per la sorgente di tracciante e per la trazione le scale sono,
rispettivamente, $E_0/L_0$, $1/T_0$, $c_{\rm ref}/T_0$ ed $E_0$.

Si decomponga il contorno meccanico in
$\partial\Omega=\overline{\Gamma_D}\cup\overline{\Gamma_N}$,
con parti relativamente aperte disgiunte e $|\Gamma_D|>0$.
Le condizioni di riferimento sono
$\mathbf d=\mathbf0$ su $\Gamma_D$ e
$\sig\mathbf n=\bar{\mathbf t}$ su $\Gamma_N$.
Su tutto il contorno si assumono impermeabilit\`a e isolamento
del tracciante,
\begin{equation}
 \mathbf q\cdot\mathbf n=0,\qquad
 \mathbf F_c\cdot\mathbf n=0
 \quad\text{su }\partial\Omega .
 \label{eq:tracerbc}
\end{equation}
La seconda condizione \`e sul flusso \emph{totale}, non soltanto
su quello diffusivo. Nel dominio anatomico si impone
$\mathbf d=0$ sull'intero contorno esterno. Le condizioni iniziali
sono
\begin{equation}
 (\mathbf d,p,c)|_{t=0}=(\mathbf d_0,p_0,c_0).
 \label{eq:initial-data}
\end{equation}

\subsection{Forma debole, struttura energetica e massa del tracciante}
\label{sec:forward-weak}

Definiamo gli spazi $V=\{\mathbf v\in H^1(\Omega)^m:
\mathbf v|_{\Gamma_D}=0\}$ e $Q=H^1(\Omega)$.
Con $(\cdot,\cdot)$ denotiamo il prodotto scalare in $L^2(\Omega)$,
con le opportune contrazioni tensoriali. Le forme bilineari sono
\[
\begin{aligned}
 a_E(\mathbf u,\mathbf v)
   &=(E(\mathbf x)\mathbb C_0:\eps(\mathbf u),\eps(\mathbf v)),\\
 a_v(\mathbf u,\mathbf v)
   &=A_v(\mathbb C_0:\eps(\mathbf u),\eps(\mathbf v)),\\
 a_h(p,w)&=(\mathbb K_h\nabla p,\nabla w),\\
 a_D(c,z)&=(\phi D\nabla c,\nabla z).
\end{aligned}
\]
Per ogni $t$ e per ogni $(\mathbf v,w,z)\in V\times Q\times Q$,
la formulazione variazionale \`e
\begin{subequations}\label{eq:weak-system}
\begin{empheq}[left=\empheqlbrace]{alignat=2}
 &a_E(\mathbf d,\mathbf v)+a_v(\CD\beta\mathbf d,\mathbf v)
       &&\notag\\[-5pt]
 &\quad -b(p,\nabla\!\cdot\mathbf v)
       &&=\langle F,\mathbf v\rangle,\label{eq:weakm}\\
 &(S_p\partial_t p+b\nabla\!\cdot\partial_t\mathbf d,w)
       &&\notag\\[-5pt]
 &\quad +a_h(p,w)
       &&=(s_f,w),\label{eq:weakp}\\
 &(\phi\CD\alpha c,z)-(c\mathbf q,\nabla z)
       &&\notag\\[-5pt]
 &\quad +a_D(c,z)
       &&=(s_c,z).\label{eq:weakc}
\end{empheq}
\end{subequations}
Il funzionale $F$ comprende $\mathbf f_s$ e la trazione
$\bar{\mathbf t}$ su $\Gamma_N$. Per definire la terza forma si
richiede, in particolare, $c\mathbf q\in L^2(\Omega)^m$.

Se $\CD\beta\mathbf d(t)\to0$ per $t\downarrow0$ quando
$0<\beta<1$, la forma meccanica impone la compatibilit\`a
\begin{equation}
 a_E(\mathbf d_0,\mathbf v)-b(p_0,\nabla\!\cdot\mathbf v)
 =\langle F(0),\mathbf v\rangle .
 \label{eq:initial-compatibility}
\end{equation}
Per $\beta=1$ compare invece
$a_v(\partial_t\mathbf d(0),\mathbf v)$; per storie frazionarie
meno regolari il limite del contributo ereditario non deve essere
presunto nullo. Nei test con carico meccanico e stato iniziale
nulli si utilizza~\eqref{eq:initial-compatibility}.

\refstepcounter{theorem}\label{prop:pv-bilanci}
\paragraph{Proposizione~\thetheorem\ (identit\`a di energia e massa).}
Si assumano $E(\mathbf x)>0$, $S_p>0$, $\mathbb K_h$ coerciva,
$A_v>0$, $0<\alpha,\beta\le1$, condizioni al contorno come sopra
e regolarit\`a sufficiente per le scelte dei test e le integrazioni
nel tempo. Posto
$\mathcal E(t)=\tfrac12a_E(\mathbf d(t),\mathbf d(t))+
\tfrac12S_p\|p(t)\|_{L^2}^2$,
vale, per ogni $t\in[0,T]$,
\begin{equation}
\begin{aligned}
 \mathcal E(t)-\mathcal E(0)
 &+\int_0^t a_h(p,p)\,\dd s\\
 &+\int_0^t a_v(\CD\beta\mathbf d,\partial_s\mathbf d)
       \,\dd s
 =W_{\rm ext}(t),
\end{aligned}
\label{eq:energy}
\end{equation}
con
$W_{\rm ext}(t)=\int_0^t[\langle F,\partial_s\mathbf d\rangle+
(s_f,p)]\,\dd s$.
I contributi cumulativi di Darcy e della memoria viscoelastica
sono non negativi. Inoltre, con
$M_\phi(t)=\int_\Omega\phi(\mathbf x)c(\mathbf x,t)\,\dd\mathbf x$
e $S_c(t)=\int_\Omega s_c(\mathbf x,t)\,\dd\mathbf x$,
\begin{equation}
 \CD\alpha M_\phi(t)=S_c(t),\qquad
 M_\phi(t)=M_\phi(0)+(I^\alpha S_c)(t),
 \label{eq:mass}
\end{equation}
dove
$(I^\alpha S_c)(t)=\Gamma(\alpha)^{-1}
\int_0^t(t-s)^{\alpha-1}S_c(s)\,\dd s$ per $0<\alpha<1$
e $I^1S_c(t)=\int_0^tS_c(s)\,\dd s$.
In assenza di sorgente, $M_\phi$ \`e costante.

\emph{Dimostrazione.}
Si usano $\mathbf v=\partial_t\mathbf d$ in
\eqref{eq:weakm} e $w=p$ in~\eqref{eq:weakp}.
La simmetria di $a_E$ e l'indipendenza temporale dei coefficienti
producono $\partial_t\mathcal E$; i due contributi misti
$\mp b(p,\nabla\!\cdot\partial_t\mathbf d)$ si cancellano.
L'integrazione da $0$ a $t$ d\`a~\eqref{eq:energy}.
Per $0<\beta<1$ la positivit\`a del contributo viscoelastico
cumulativo segue dalla rappresentazione a misura positiva
$k_\beta(t)=(\sin\pi\beta/\pi)
\int_0^\infty r^{\beta-1}e^{-rt}\,\dd r$
applicata a $\eps(\partial_t\mathbf d)$, come esplicitato
nell'Appendice~\ref{app:mms}; per $\beta=1$ esso \`e
$\int_0^t a_v(\partial_s\mathbf d,\partial_s\mathbf d)\,\dd s$.
La forma $a_h$ \`e non negativa per la positivit\`a di
$\mathbb K_h$. Infine, scegliendo $z=1\in Q$ in
\eqref{eq:weakc}, i termini di avvezione e diffusione si annullano;
l'indipendenza temporale di $\phi$ permette di portare l'integrale
spaziale all'interno della derivata di Caputo. L'identit\`a
$\CD\alpha M_\phi=S_c$ si integra applicando $I^\alpha$.
\hfill$\square$

L'identit\`a~\eqref{eq:energy} \`e un bilancio \emph{continuo} e
non implica automaticamente dissipazione numerica monotona a ogni
passo dell'algoritmo temporale. Le corrispondenti identit\`a discrete
e la verifica della massa pesata con $\phi$ sono documentate
nella Sezione~\ref{sec:numerics} e nell'Appendice~\ref{app:mms}.
